\documentclass[11pt,a4paper]{article}
\usepackage[margin=0.5in, top=60pt]{geometry}
\usepackage{graphicx}
\graphicspath{ {./} }
\usepackage{multirow}
\usepackage{makecell}
\usepackage[table,xcdraw]{xcolor}
\usepackage{setspace}
\usepackage{pgffor}
\usepackage{titlesec}

\newcommand{\code}[1]{
    \texttt{#1}
}

\setlength\parindent{0pt}
\setlength\parskip{8pt}
\titlespacing*{\section} {0pt}{0ex}{0ex}



\usepackage{graphicx}
\usepackage{fancyhdr}
\fancyhf{}
\setlength\headheight{20pt}
\fancyhead[L]{\includegraphics[width=1cm]{images/diu.png} \ \includegraphics[width=1.4cm]{images/icpc.png}}
\fancyhead[C]{ICPC Dhaka Regional Contest, 2024 Editorial\\ Hosted By: Daffodil International University}
\fancyhead[R]{\includegraphics[width=1.5cm]{images/baps.png}}
\renewcommand\headrulewidth{1pt}
\pagestyle{fancy}


\begin{document}

\begin{titlepage}
    \begin{center}
        \textbf{\large International Collegiate Programming Contest (ICPC)\\
        2024 Asia Dhaka Regional Contest\\
        Onsite Round - Editorial\\
        }
        \vspace{3cm}
        \textbf{\large Hosted By \\
        Department of Computer Science and Engineering\\
        Daffodil International University\\}
        \vspace{7cm}
        \textbf{\large Supported by\\}
        \vspace{1cm}
        \includegraphics[width=3cm]{images/diu.png} \ \ 
        \includegraphics[width=4cm]{images/icpc.png} \ \ 
        \includegraphics[width=4cm]{images/baps.png} \\
        \vspace{1cm}
        \includegraphics[width=8cm]{images/dbbl.png} \ \ 
        \includegraphics[width=4cm]{images/regional.png} \ \ 
        
    \end{center}
\end{titlepage}

\section*{Problem A. Are the Nodes Reachable?}
Problem Setter: Raihat Zaman Neloy\\
Tester: Rumman Mahmud, Pritom Kundu, Nafis Sadique\\
Cetegory: Ad-hoc\\
Total Solved: \\
First to Solve: \\
\\
We can split the problem is two subproblem.
\paragraph*{When $ans < 32$.}
 How do we solve it? First for each vertex calculate the list of vertex reachable from it. Also 
calculate the list of vertex from which it is reachable. Both can be done in $O(\frac{max(n) \cdot m}{64})$ using bitsets.
Now assume that vertex $a$ is reachable from vertex $b$. In that case vertices $b-1$ or $b+1$ will be reachable by 
adding an edge of cost 1. We can find all the vertices from which vertex $a$ is reachable with 1 cost by simply
binary shifting the bitset by 1 in both direction and apply the bitwise OR operator on all three (no change, left shift, right shift).
If we do this again then we find the vertices 2 distance away and so on. We keep doing this for distances upto 31. The
complexity to do this is $O(\frac{n^2*32}{64})$. 
\\
\\
Now, lets split the bitsets into $64$ sized buckets. Technically we could've used an unsigned long long array to perform
the bitset operations in the first place. For each vertex we will have 33 bitset buckets. 32 of them are for
the shifted bitsets from upstream and one for the downstream.
Then for the query, we check if the first upstream bucket of V with 31 shift and the downstream bucket of U has any match. We 
can do that by simply performing a bitwise AND operation. If there is a match then we check 30 shift and keep reducing it. If 
no match is found then we move to the next bucket. At the end in this approach we will need $O(\frac{n}{64}+32)$ operations
per query.
\paragraph*{When $ans \ge 32$.}
We again keep a bucket of vertices reachable from a vertex. Vertices numbered from [1..64] are in one bucket, [65, 128] are 
in second buckets and so on. For each bucket we only keep the minimum id of the vertex reachable from it. Similarly we do it 
for the upstream vertices. For a query, we merge those two list of vertices (min and max of each bucket), sort them and find the
smallest gap in consecutive vertices (one must be upstream while the other must be downstream). This takes $O(\frac{n*2}{64})$ per
query.
\\ 
\\
Alternatively, we can do the same using 128 bit integers. The performance is comparable.

\section*{Problem B. Yet Another Crossover Episode}
Problem Setter: Shahjalal Shohag \\
Tester: Rumman Mahmud, Pritom Kundu, Jubayer Rahman, Nafis Sadique \\
Cetegory: Dynamic Programming, Bit Manipulation \\
Total Solved:  \\
First to Solve: \\

\section*{Problem C. Cut the Stick, Share You Must}
Problem Setter: Rumman Mahmud \\
Tester: Shahjalal Shohag, Pritom Kundu, Nafis Sadique \\
Cetegory: Math, Number Theory \\
Total Solved:  \\
First to Solve: \\

\section*{Problem D. CatGPT}
Problem Setter: Anik Sarker \\
Tester: Rumman Mahmud, Pritom Kundu, Nafis Sadique \\
Cetegory: Graph, Offline Processing \\
Total Solved:  \\
First to Solve: \\
\\
There can be only 26 clowders. We can use this information. First, lets process
all the event in the ascending order of $R_i$. We then process the cats from left
to right. After processing each cat we keep in memory the most recent position we 
saw a cat of clowder $C_i$. Then we process the events ended at the current position.
Because we know the last position for the clowder, we know that event will merge
clowders that have their $last\ position > L_i$. We keep the list of merged clowders
for the event in memory as well.\\
\\
For the queries, we also process them in the ascending order of $y_i$. But this time we 
iterate over the events from left to right. Since we know for an event, the clowders that
will be merged, we can generate a list of unordered $(i,j)$ pairs that will be merged for
an event. There can't be more than $\frac{25 \cdot 26}{2}$ pairs. We keep track of the last event
they were merged, similar to the previous section. Finally for an event, we process the 
queries that ended in that event. From the list of the last position of the pairs, we
can have a graph of 26 vertices. We merge vertices using union-find. Then we check the largest
merged clowder and output their size.

\section*{Problem E: Quasi-binary Representations}
Problem Setter: Pritom Kundu \\
Tester: Shahjalal Shohag, Pritom Kundu, Jubayer Rahman, Nafis Sadique \\
Cetegory: Dynamic Programming, Math \\
Total Solved:  \\
First to Solve: \\

\section*{Problem F: Flowers}
Problem Setter: Hasinur Rahman \\
Tester: Nafis Sadique \\
Cetegory: Hashing, Divide \& Conquer \\
Total Solved:  \\
First to Solve: \\

\section*{Problem G: Library Function vs Keyword}
Problem Setter: Shahriar Manzur \\
Tester: Jubayer Rahman, Nafis Sadique, Rumman Mahmud \\
Cetegory: Ad-hoc \\
Total Solved:  \\
First to Solve: \\
\\
For $strlen(line)$, we simply need to find the position of the first null(`\textbackslash 0') character.
The number of characters before that is the answer. If no null characters exist then
the answer would be the string length.
\\ \\
For $sizeof(line)$, cound the number of null in the string. Then the answer is 
$string\ length - number\ of\ null\ characters + 1$.


\section*{Problem H: Hand Cricket}
Problem Setter: Kazi Md Irshad \\
Tester: Pritom Kundu, Jubayer Rahman, Nafis Sadique \\
Cetegory: Data Structure, Math, Probability \\
Total Solved:  \\
First to Solve: \\

\section*{Problem I: In Search of a Kind Person}
Problem Setter: Md. Imran Bin Azad \\
Tester: Shahjalal Shohag, Rumman Mahmud, Raihat Zaman Neloy, Muhiminul Islam Osim \\
Cetegory: Ad-hoc \\
Total Solved:  \\
First to Solve: \\

\section*{Problem J: The Taxman}
Problem Setter: Aminul Haque \\
Tester: Rumman Mahmud,  Muhiminul Islam Osim \\
Cetegory: Binary Search, Math \\
Total Solved:  \\
First to Solve: \\

\section*{Problem K: Packet Transmission}
Problem Setter: Ashraful Islam \\
Tester: Pritom Kundu, Nafis Sadique \\
Cetegory: Least Common Ancestor \\
Total Solved:  \\
First to Solve: \\
\\
Every query we receive can be classified in the following way.
\begin{itemize}
    \item The query packets (source, destination pair) share a common path in the tree. 
    \begin{itemize}
        \item The packets are going in the same direction. In that case we need to assume one of the
        packets will never wait. So, the other packet will arrive at the first vertex of the common path,
        wait for the first packet to arrive and then go after it. Since an edge can't be used by
        multiple packets at the same time, the other packet will have to wait until the first packet 
        crossed that edge. Keep doing that until they reach the end of common path and then they can 
        go their own ways. If we think carefully, we can see that the other packet will have to wait
        an additional time totalling the maximum edge cost on the common path. However if the first
        packet comes early then the wait time reduces. Do the same the other way around and take
        the minimum of the two.
        \item The packets are going in the different directions. So, we need both packets to start
        their journey. Eventually they might meet on the two sides of an edge where if one starts crossing,
        the other must wait. Make one of them wait and the other one go and calculate the time it would 
        take for both them to reach their destinations. Do it both ways and take the minimum time. However if
        they don't meet at the edge then they don't have to stop and their actual time to reach both 
        their destinations is the answer. 
    \end{itemize}
    \item The query packets use completely different paths. In that case the time it takes for both to reach
    their destination is the answer.
\end{itemize}
We can calculate all of these using Sparse Table to calculate Least Common Ancestors in a tree. Alternatively 
we can use heavy-light decomposition, but that may be very slow. There are lot of corner cases, so careful 
implementation is necessary.

\section*{Problem L: Unhappy Team}
Problem Setter: Nafis Sadique \\
Tester: Shahjalal Shohag, Rumman Mahmud, Jubayer Rahman \\
Cetegory: Dynamic Programming \\
Total Solved:  \\
First to Solve: \\
\\ 
The trick is to select the unhappiness score of someone and count how many times that score appear
in the top K values. We can keep a DP state $(bit\_mask, bigger\_scores\_count, selected\_score\_appeared)$.
Basically we will put people one after another, everytime calculating their unhappiness score. If the 
score is smaller than the selected score then we ignore it. If the score is bigger then $bigger\_scores\_count$
is incremented by one. If the score is the same is the selected score then we can consider it as bigger or 
mark $selected\_score\_appeared$ as $true$. This DP will have the complexity of $O(2^n\cdot n^3)$. With some
other minor optimizations, this runs reasonably fast.

\end{document}